\chapter{KAJIAN PUSTAKA}
\label{sec:chap2_kajian_pustaka}

\section{Hasil Penelitian Terdahulu}
Penelitian ini didukung oleh berbagai acuan dari penelitian sebelumnya yang relevan dengan pengembangan drone otonom dan sistem deteksi objek untuk misi penyelamatan. Tinjauan ini berfungsi sebagai dasar perbandingan dan penentu posisi penelitian dalam perkembangan teknologi terkini.

\subsection{Autonomous Search and Rescue Drone}
Penelitian oleh Elashaal dkk. (2024) mengembangkan sistem drone otonom SAR menggunakan platform Holybro S500 dengan \textit{flight controller} Pixhawk 2.4.8 dan Raspberry Pi 4. Sistem ini menerapkan algoritma YOLOv4-tiny untuk deteksi manusia secara \textit{real-time}. Hasilnya menunjukkan akurasi mAP sebesar 96,44\% pada citra \textit{top-view}, namun memiliki kendala pada limitasi komputasi Raspberry Pi 4 yang menyebabkan latensi tinggi dan \textit{frame rate} yang rendah dalam kondisi medan kompleks. Selain itu, pengujian lapangan masih terbatas pada area spesifik seperti gurun dan pantai \cite{Elashaal2024ASARD}.

\subsection{Real-Time Human Detection and Navigation for SAR}
Jayalath dan Munasinghe (2021) mengintegrasikan deteksi manusia berbasis \textit{TensorFlow Fast R-CNN} dengan navigasi otonom pada Raspberry Pi 4 menggunakan \textit{dataset} Okutama-Action. Drone mampu menghitung koordinat GPS manusia berdasarkan posisi \textit{bounding box}. Meskipun akurasi mencapai 84,1\%, waktu pemrosesan yang dibutuhkan mencapai 8 detik per \textit{frame}, yang dinilai kurang responsif untuk misi otonom dinamis. Selain itu, sistem sangat bergantung pada kondisi pencahayaan alami \cite{Jayalath2021DroneHumanID}.

\subsection{Target Detection, Tracking, and Avoidance System for UAVs}
Varatharasan dkk. (2019) berfokus pada efisiensi perangkat keras berbiaya rendah dengan menggunakan Raspberry Pi 3B+ dan Intel Movidius NCS2. Penelitian ini menonjolkan penggunaan \textit{MAVSDK API} untuk kendali mode \textit{Offboard}. Kelebihannya adalah stabilitas sistem pada perangkat berdaya rendah, namun memiliki keterbatasan pada ketepatan estimasi kedalaman yang masih bergantung pada pendekatan \textit{stadiametric rangefinding} \cite{Varatharasan2019TargetDetection}.

\subsection{Posisi Penelitian terhadap Penelitian Terdahulu}
Berdasarkan tinjauan terhadap ketiga penelitian terdahulu di atas, terdapat beberapa keterbatasan operasional yang menjadi dasar pengembangan penelitian ini. Pertama, penelitian terdahulu menggunakan papan komputasi tunggal generasi lama yang memiliki keterbatasan komputasi untuk mengeksekusi model deteksi secara responsif. Kedua, arsitektur perangkat lunak yang diusulkan belum sepenuhnya mengintegrasikan pemisahan tugas secara \textit{multi-thread} atau \textit{asynchronous} untuk menjaga stabilitas navigasi dan inferensi secara simultan. 

Penelitian ini bertujuan untuk mengisi keterbatasan tersebut dengan mengintegrasikan Raspberry Pi 5 sebagai \textit{companion computer} untuk pemrosesan lokal, menerapkan optimasi model YOLOv8n kelas tunggal manusia pada format ONNX, serta merancang arsitektur perangkat lunak terdistribusi untuk memisahkan beban komputasi navigasi otonom dan pengolahan visual.

\section{Dasar Teori}
Subbab ini memaparkan teori-teori dasar yang menjadi landasan teknis dalam pengerjaan penelitian sistem drone otonom.

\subsection{Unmanned Aerial Vehicle (UAV)}
\textit{Unmanned Aerial Vehicle} (UAV) atau drone adalah pesawat tanpa awak yang beroperasi secara mandiri melalui sensor tersemat atau kendali jarak jauh. UAV modern mengintegrasikan sensor \textit{Inertial Measurement Unit} (IMU), GPS, dan algoritma kontrol otonom. Pada penelitian ini digunakan tipe \textit{rotary-wing quadcopter} karena kemampuan \textit{hover} (diam di udara) dan fleksibilitas manuver yang stabil untuk operasi pengamatan visual di area bencana \cite{Austin2021}. Konfigurasi \textit{quadcopter} menggunakan empat rotor yang diatur dalam formasi silang (\textit{X-frame}) memberikan redundansi kontrol yang memadai untuk menjaga kestabilan penerbangan saat membawa beban komputasi tambahan \cite{Quan2017MulticopterDesign}.

\subsection{Holybro S500}
Holybro S500 adalah platform \textit{quadcopter} riset berbasis \textit{X-frame} yang sering digunakan dalam pengembangan UAV akademik. Frame ini terbuat dari serat karbon dan plastik ABS, mendukung kapasitas \textit{payload} hingga 1,5 kg. S500 dinilai memadai untuk integrasi \textit{companion computer} dan algoritma kecerdasan buatan karena stabilitas strukturalnya yang baik \cite{Holybro2023}. Platform ini menyediakan \textit{Power Distribution Board} (PDB) terintegrasi yang memudahkan distribusi daya ke seluruh komponen motor \textit{brushless} dan \textit{Electronic Speed Controller} (ESC).

\subsection{Flight Controller Pixhawk dan Firmware PX4}
\textit{Flight controller} berfungsi sebagai unit kendali terbang utama yang memproses data sensor dan menghasilkan sinyal kontrol untuk motor penggerak. Pada penelitian ini digunakan Pixhawk 6C yang merupakan generasi terbaru dari keluarga Pixhawk dengan prosesor STM32H743 \cite{Pixhawk6C2024}. PX4 adalah perangkat lunak \textit{autopilot open-source} yang mengelola \textit{flight stack} penerbangan and menyediakan estimasi posisi melalui algoritma \textit{Extended Kalman Filter} (EKF2) berbasis \textit{sensor fusion} dari data IMU, barometer, dan GPS \cite{Meier2015PX4, Faragher2012EKF}. PX4 mendukung mode penerbangan \textit{Offboard} yang memungkinkan \textit{companion computer} mengirimkan perintah navigasi berbasis kecepatan (\textit{velocity setpoints}) secara langsung melalui protokol MAVLink \cite{PX4DevGuide2024}. Selain itu, PX4 menyediakan mekanisme keselamatan terintegrasi berupa \textit{Battery Return to Launch} (Battery RTL) untuk mengembalikan drone secara otonom ketika kapasitas daya baterai menurun ke ambang batas kritis.

\subsection{Raspberry Pi 5 sebagai Companion Computer}
Raspberry Pi 5 merupakan \textit{companion computer} generasi terbaru yang membawa peningkatan performa signifikan dengan prosesor ARM Cortex-A76 quad-core 2,4 GHz dan GPU VideoCore VII. Dengan RAM hingga 8 GB, perangkat ini mampu menjalankan beban kerja \textit{Edge AI} seperti inferensi \textit{deep learning} secara lokal \cite{RaspberryPi2024}. Karakteristik dayanya yang sensitif (5V 5A) memerlukan sistem regulasi tegangan yang stabil untuk menghindari kegagalan \textit{booting} atau \textit{voltage sag} saat beban komputasi puncak. Dalam konteks drone otonom, Raspberry Pi 5 diposisikan sebagai pusat pemrosesan yang menjalankan seluruh arsitektur \textit{multi-thread}, mengeksekusi model deteksi manusia secara \textit{real-time}, dan melayani transmisi data ke antarmuka pemantauan.

\subsection{Komputasi Tepi (Edge Computing)}
\textit{Edge computing} adalah paradigma komputasi yang memindahkan proses pengolahan data mendekati sumber data alih-alih mengirimkannya ke pusat data (\textit{cloud}) yang jauh \cite{Shi2016EdgeComputing}. Dalam konteks UAV, pendekatan ini dinilai penting karena misi pencarian korban seringkali dilakukan di area terpencil dengan keterbatasan \textit{bandwidth} jaringan. Dengan menempatkan inferensi model \textit{deep learning} langsung pada \textit{companion computer} di atas drone, latensi pengambilan keputusan dapat diminimalkan dibandingkan arsitektur berbasis \textit{cloud} \cite{Chen2019DeepEdge, Mao2021EdgeUAV}. Implementasi \textit{edge computing} pada penelitian ini mengharuskan optimasi pada beberapa aspek, meliputi mode sistem operasi (\textit{headless}), resolusi input model, dan format inferensi yang efisien.

\subsection{Deep Learning dan Convolutional Neural Network}
\textit{Deep learning} merupakan cabang dari \textit{machine learning} yang menggunakan jaringan saraf tiruan berlapis dalam (\textit{deep neural networks}) untuk mengekstraksi representasi fitur hierarkis dari data mentah secara otomatis \cite{LeCun2015DeepLearning}. \textit{Convolutional Neural Network} (CNN) adalah arsitektur \textit{deep learning} yang dirancang khusus untuk memproses data berbentuk grid seperti citra digital. CNN menggunakan operasi konvolusi untuk mengekstraksi fitur spasial, diikuti oleh lapisan \textit{pooling} untuk reduksi dimensi, dan lapisan \textit{fully connected} untuk klasifikasi akhir. Arsitektur ini menjadi fondasi utama pada algoritma deteksi objek modern termasuk keluarga YOLO yang digunakan dalam penelitian ini.

\subsection{YOLO (You Only Look Once) v8}
YOLOv8 adalah algoritma deteksi objek berbasis CNN yang melakukan deteksi dan klasifikasi dalam satu langkah (\textit{single-stage detector}). Berbeda dengan detektor dua tahap (\textit{two-stage detector}) seperti Fast R-CNN yang memerlukan tahap proposal region terpisah, arsitektur \textit{single-stage} pada YOLO memproses seluruh citra dalam satu kali \textit{forward pass} sehingga menghasilkan latensi yang lebih rendah \cite{Terven2023YOLOv8Review}. Versi kedelapan ini memperkenalkan arsitektur \textit{anchor-free} dan modul C2f (\textit{Cross Stage Partial with two convolutions}) yang meningkatkan akurasi. Varian YOLOv8n (\textit{nano}) memiliki jumlah parameter terkecil dalam keluarga YOLOv8 dengan ukuran model sekitar 6 MB, sehingga dapat digunakan untuk \textit{deployment} pada perangkat \textit{edge} dengan mempertahankan laju \textit{frame} yang memadai \cite{Zhou2024SODYOLOv8}.

\subsection{Transfer Learning dan Fine-Tuning}
\textit{Transfer learning} adalah teknik pembelajaran mesin yang memanfaatkan pengetahuan (\textit{knowledge}) yang telah dipelajari oleh model pada suatu domain untuk diterapkan pada domain target yang berbeda \cite{Zhuang2020TransferLearning}. Dalam konteks deteksi objek, pendekatan \textit{fine-tuning} dilakukan dengan mengambil model yang telah dilatih pada \textit{dataset} besar (misalnya COCO) kemudian melatih ulang lapisan tertentu menggunakan \textit{dataset} spesifik sesuai kebutuhan aplikasi. Pada penelitian ini, model YOLOv8n yang telah di-\textit{pretrain} pada 80 kelas COCO dilakukan \textit{fine-tuning} menjadi detektor kelas tunggal khusus manusia (\textit{person-only}). Strategi ini memungkinkan model memfokuskan kapasitas ekstraksi fitur pada kontur dan siluet manusia, sehingga meningkatkan akurasi deteksi sekaligus mengurangi beban komputasi untuk klasifikasi objek non-target.

\subsection{Dataset COCO (Common Objects in Context)}
COCO (\textit{Common Objects in Context}) merupakan \textit{dataset} berskala besar yang digunakan sebagai \textit{benchmark} untuk evaluasi algoritma deteksi objek, segmentasi, dan \textit{captioning} \cite{Lin2014COCO}. Dataset COCO2017 memiliki lebih dari 200.000 citra berlabel dengan 80 kategori objek. Pada penelitian ini, dataset COCO2017 difilter untuk mengambil hanya anotasi kelas \textit{person} dengan total 244.799 instans \textit{bounding box} target. Penggunaan \textit{dataset} ini memastikan model memiliki variasi data pelatihan yang mencakup perbedaan pose tubuh, pencahayaan, resolusi, serta kondisi oklusi parsial.

\subsection{Format ONNX dan ONNX Runtime}
ONNX (\textit{Open Neural Network Exchange}) adalah format standar terbuka untuk merepresentasikan model \textit{machine learning} yang memungkinkan interoperabilitas antar kerangka kerja (\textit{framework}) seperti PyTorch, TensorFlow, dan lainnya \cite{Bai2019ONNX}. ONNX Runtime merupakan mesin inferensi yang dioptimasi untuk mengeksekusi model berformat ONNX pada berbagai platform perangkat keras \cite{ONNXRuntime2019}. Pada penelitian ini, model YOLOv8n hasil \textit{fine-tuning} dalam format PyTorch (\texttt{best.pt}) diekspor ke format ONNX dengan resolusi input 320$\times$320 piksel untuk diimplementasikan pada Raspberry Pi 5. Konversi ke format ONNX memberikan keuntungan berupa optimasi graf komputasi dan kompatibilitas lintas platform.

\subsection{Metrik Evaluasi Deteksi Objek}
Evaluasi performa model deteksi objek pada penelitian ini menggunakan beberapa metrik standar yang digunakan dalam literatur \textit{computer vision} \cite{Padilla2021mAPSurvey}:
\begin{itemize}
	\item \textit{Precision}: rasio prediksi benar positif terhadap seluruh prediksi positif.
	\item \textit{Recall}: rasio prediksi benar positif terhadap seluruh objek target yang sebenarnya.
	\item \textit{F1-Score}: rata-rata harmonik dari \textit{precision} dan \textit{recall}.
	\item mAP@50 (\textit{mean Average Precision} pada IoU 0,5): rata-rata \textit{average precision} pada ambang batas \textit{Intersection over Union} sebesar 50\%.
	\item mAP@50-95: rata-rata mAP pada rentang IoU dari 0,5 hingga 0,95 dengan interval 0,05.
\end{itemize}

\subsection{MAVLink dan MAVSDK}
\textit{MAVLink} (\textit{Micro Air Vehicle Link}) adalah protokol komunikasi ringan berbasis paket biner untuk pertukaran data antara \textit{flight controller} dan perangkat eksternal seperti \textit{companion computer} atau \textit{Ground Control Station} (GCS) \cite{MAVLink2024}. Protokol ini dirancang dengan \textit{overhead} minimal untuk mendukung komunikasi pada sistem tertanam. Sementara itu, \textit{MAVSDK} adalah \textit{software development kit} modern yang menyediakan API (C++, Python, Swift) untuk mengakses fungsi navigasi tingkat tinggi seperti \textit{offboard control}, telemetri, dan perintah misi \cite{mavsdk_doc2025}. Pada penelitian ini, MAVSDK-Python digunakan untuk mengirimkan perintah kecepatan (\texttt{VelocityNedYaw}) ke \textit{flight controller} dan membaca data telemetri penerbangan.

\subsection{Sensor Visual (Webcam USB)}
Kamera \textit{webcam} dengan antarmuka USB digunakan sebagai sensor visual utama karena kemudahan integrasi (\textit{plug-and-play}) pada ekosistem Linux Raspberry Pi melalui subsistem \textit{Video4Linux2} (V4L2) \cite{V4L2_2024}. Pengaturan parameter kamera dilakukan melalui utilitas \texttt{v4l2-ctl} untuk menyediakan aliran video yang mendukung kebutuhan pemrosesan visi komputer \cite{Bradski2000OpenCV}. Pada penelitian ini digunakan Logitech C922 yang mampu menangkap aliran video dengan tingkat latensi rendah melalui antarmuka USB 3.0.

\subsection{GStreamer}
GStreamer adalah \textit{framework} multimedia \textit{open-source} yang menyediakan arsitektur berbasis \textit{pipeline} untuk mengolah aliran data audio dan video \cite{GStreamer2024}. Arsitektur \textit{pipeline} pada GStreamer menggunakan elemen-elemen modular yang dihubungkan secara serial maupun paralel. Pada penelitian ini, GStreamer digunakan sebagai subsistem akuisisi video yang menangkap aliran data dari kamera USB melalui elemen \texttt{v4l2src}, kemudian mendistribusikan \textit{frame} melalui mekanisme percabangan (\textit{tee pipeline}) ke dua jalur: jalur \textit{streaming} untuk pemantauan operator dan jalur \textit{appsink} untuk inferensi model YOLOv8.

\subsection{OpenCV}
\textit{OpenCV} (\textit{Open Source Computer Vision Library}) merupakan pustaka pemrosesan citra dan visi komputer yang digunakan dalam pengembangan sistem berbasis visual \cite{Bradski2000OpenCV, song2022_performance_opencv}. Dalam penelitian ini, OpenCV berperan dalam \textit{pipeline} pra-pemrosesan (perubahan ukuran \textit{frame}, konversi format warna, normalisasi nilai piksel) dan pasca-pemrosesan (visualisasi kotak pembatas hasil deteksi dan pengukuran FPS). OpenCV juga menyediakan antarmuka untuk konversi format data \textit{buffer} dari GStreamer menjadi representasi matriks NumPy yang dapat diproses oleh model deteksi.

\subsection{Finite State Machine (FSM)}
\textit{Finite State Machine} (FSM) atau mesin keadaan berhingga adalah model komputasi yang terdiri dari sekumpulan \textit{state} (keadaan), transisi antar \textit{state}, dan aksi yang dieksekusi berdasarkan kondisi tertentu \cite{Wagner2006FSM}. Pada sistem drone otonom, pendekatan FSM digunakan untuk merepresentasikan perilaku navigasi drone selama misi pencarian korban. Setiap \textit{state} merepresentasikan fase operasi tertentu seperti pemindaian area (\texttt{SCAN}), pendekatan target (\texttt{APPROACH}), dokumentasi manusia yang terdeteksi (\texttt{DOCUMENT}), dan pergeseran posisi lateral (\texttt{DISPLACING}). Transisi antar \textit{state} dipicu oleh data telemetri dan hasil deteksi visual, dengan pengecualian prosedur \textit{Return to Launch} (RTL) yang dapat diinisiasi oleh sistem saat baterai mencapai level kritis sesuai konfigurasi PX4.

\subsection{Exponential Moving Average (EMA)}
\textit{Exponential Moving Average} (EMA) adalah metode pemulusan (\textit{smoothing}) data deret waktu yang memberikan bobot eksponensial pada pengamatan terbaru \cite{Hunter1986EMA}. Formula EMA didefinisikan sebagai:
\begin{equation}
	S_t = \alpha \cdot x_t + (1 - \alpha) \cdot S_{t-1}
\end{equation}
di mana $S_t$ adalah nilai EMA pada waktu $t$, $x_t$ adalah pengamatan terbaru, dan $\alpha$ ($0 < \alpha \leq 1$) adalah faktor pemulusan. Pada sistem drone otonom, EMA diterapkan untuk menghaluskan koordinat posisi target (\textit{center x} dan \textit{center y}) dari hasil deteksi visual sehingga perintah navigasi yang dihasilkan lebih stabil terhadap fluktuasi posisi \textit{bounding box} antar \textit{frame}.

\subsection{Formula Jarak Haversine}
Formula \textit{Haversine} adalah persamaan navigasi matematis yang digunakan untuk menghitung jarak lingkaran besar (\textit{great-circle distance}) antara dua titik pada permukaan bola bumi berdasarkan koordinat garis lintang dan garis bujur \cite{Sinnott1984Haversine}. Formula ini digunakan untuk menghitung jarak spasial antar objek. Persamaan \textit{Haversine} didefinisikan sebagai:
\begin{equation}
	a = \sin^2\left(\frac{\Delta \phi}{2}\right) + \cos(\phi_1) \cdot \cos(\phi_2) \cdot \sin^2\left(\frac{\Delta \lambda}{2}\right)
\end{equation}
\begin{equation}
	d = 2 \cdot R \cdot \operatorname{atan2}\left(\sqrt{a}, \sqrt{1-a}\right)
\end{equation}
di mana $\phi$ adalah garis lintang, $\lambda$ adalah garis bujur, $\Delta \phi$ dan $\Delta \lambda$ adalah selisih koordinat antar kedua titik, serta $R$ adalah radius rata-rata bumi (sekitar 6.371 km). Dalam konteks sistem drone otonom ini, formula \textit{Haversine} diimplementasikan untuk menciptakan radius eksklusi (\textit{exclusion radius}). Komputasi ini berfungsi menolak target manusia baru yang titik koordinat GPS-nya terlalu dekat dengan manusia yang telah didokumentasikan sebelumnya, mencegah duplikasi penargetan selama proses \textit{scouting}.

\subsection{Arsitektur Multi-Thread dan Asynchronous}
Arsitektur \textit{multi-thread} memungkinkan sebuah program menjalankan beberapa alur eksekusi secara bersamaan dalam satu proses, sehingga tugas-tugas yang bersifat independen dapat diproses secara paralel \cite{Tanenbaum2023OperatingSystems}. Pada sistem drone otonom yang memerlukan pemrosesan inferensi AI, komunikasi telemetri, dan \textit{streaming} video secara simultan, pemisahan proses ke dalam \textit{thread} independen diperlukan untuk mencegah hambatan (\textit{bottleneck}) komputasi. Pemrograman \textit{asynchronous} menggunakan pustaka \texttt{asyncio} pada Python memungkinkan operasi input/output (\textit{I/O-bound}) seperti komunikasi MAVLink berjalan secara non-\textit{blocking} dengan frekuensi tinggi. Kombinasi kedua pendekatan ini membantu memastikan bahwa latensi komputasi pada satu subsistem tidak menghentikan pengiriman sinyal kontrol ke drone.

\subsection{Sistem Manajemen Daya (UBEC dan Baterai LiPo)}
Dalam sistem drone otonom, stabilitas daya untuk unit komputasi adalah aspek penting. Baterai \textit{Lithium Polymer} (LiPo) digunakan sebagai sumber energi utama karena memiliki rasio energi terhadap berat yang tinggi dan kemampuan pelepasan arus (\textit{C-Rating}) yang besar \cite{Traub2011BatteryPowered}. Modul \textit{Universal Battery Elimination Circuit} (UBEC) 5A berfungsi sebagai regulator tegangan \textit{step-down} yang menurunkan tegangan baterai LiPo 4S (nominal 14,8V) menjadi 5V stabil untuk Raspberry Pi 5 \cite{mcnulty2022_liion_review}. Pemilihan UBEC dibandingkan modul regulator konvensional didasarkan pada kebutuhan respons transien yang cepat dan kemampuan penanganan arus puncak saat inferensi AI berjalan bersamaan dengan operasi penerbangan.

\subsection{Manufaktur Aditif (Cetak 3D FDM)}
\textit{Fused Deposition Modeling} (FDM) adalah metode manufaktur aditif yang membentuk objek tiga dimensi melalui proses deposisi material termoplastik secara lapis per lapis \cite{Ngo2018AdditiveManufacturing}. Material \textit{Polylactic Acid} (PLA) dipilih karena memiliki sifat mekanis yang memadai, kemudahan pencetakan, dan berat jenis yang rendah sehingga cocok untuk komponen struktural drone \cite{Farah2016PLA}. Pada penelitian ini, seluruh komponen mekanik pendukung seperti dudukan kompartemen, pelindung \textit{companion computer}, dan \textit{bracket} kamera dirancang menggunakan pemodelan parametrik pada perangkat lunak CAD dan dicetak menggunakan metode FDM.

\subsection{Arsitektur Backend (Flask dan SocketIO)}
Flask adalah \textit{micro web framework} berbasis Python yang menyediakan infrastruktur untuk pengembangan aplikasi web dan \textit{REST API} \cite{Flask2024}. SocketIO merupakan pustaka komunikasi dua arah (\textit{bidirectional}) berbasis protokol \textit{WebSocket} yang memungkinkan pertukaran data antara server dan klien \cite{SocketIO2024, Fette2011WebSocket}. Pada sistem drone otonom, Flask bertindak sebagai pusat orkestrasi \textit{backend} yang mengelola seluruh modul perangkat lunak, sementara SocketIO digunakan untuk memancarkan (\textit{broadcast}) data telemetri dan hasil deteksi ke antarmuka pemantauan operator.

\subsection{Antarmuka Frontend (React)}
React adalah pustaka JavaScript yang dikembangkan oleh Meta untuk membangun antarmuka pengguna (\textit{user interface}) berbasis komponen reaktif \cite{React2024}. Arsitektur berbasis komponen pada React memungkinkan pembaruan tampilan secara efisien melalui mekanisme \textit{virtual DOM}, sehingga hanya elemen yang mengalami perubahan data yang dirender ulang. Pada penelitian ini, React digunakan untuk membangun antarmuka dasbor pemantauan (\textit{monitoring dashboard}) yang menampilkan metrik penerbangan, siaran video kamera, dan status \textit{state machine} otonom secara \textit{real-time}.

\subsection{Remote Controller dan Safety Override}
Sistem \textit{Remote Controller} (RC) berfungsi sebagai instrumen kendali manual dan protokol keselamatan utama. Melalui frekuensi radio (2,4 GHz) dan protokol SBUS, RC memungkinkan operator melakukan \textit{override} atau mengambil alih kontrol drone secara langsung apabila sistem otonom mengalami anomali atau terjadi kondisi darurat di udara \cite{Austin2021}. Mekanisme \textit{override} ini bertujuan untuk memastikan bahwa pilot manusia memiliki otoritas tertinggi atas penerbangan drone, terlepas dari kondisi sistem otonom yang sedang berjalan.